% !TEX root = main.tex
\section{관련연구와 설계의 위치}
본 절은 DPP 주장 검증, 비공개 이력 복원, 사용 제한의 대상과
집행 방식을 구분하여 비교한다. 문헌에서 특정 절차를 확인하
지 못했다는 사실을 그 방식으로 해당 절차를 구현할 수 없다는
근거로 사용하지 않는다.

\subsection{DPP 검증과 공급망 생애주기}
Udokwu와 Craß는 지속가능성 주장과 이를 뒷받침하는 비공개
자료를 분리하고 영지식증명으로 주장 조건을 검증하는 DPP
구성을 제안한다 [3]. Babel 등은 생산 단계의 검증 가능한 자
격증명을 비공개 토큰의 이전과 결합하고, 사용한 지분을 소진
하여 반복 사용을 방지하는 생애주기를 다룬다 [2]. 따라서 기
존 DPP 연구를 모두 독립 문서의 검증에 한정된 것으로 정리하
는 것은 적절하지 않다. 본 연구는 이와 같은 단계별 검증에 사
후에 지정된 문제 시작점의 후속 객체 식별과 차단을 연결하는
데 초점을 둔다. 회계 크레딧 배분과 디지털 소비 방지가 실제
재료의 구성이나 동일 물량의 외부 중복 등록까지 증명하는 것
은 아니다.

Sahai 등의 공급망 시스템은 문서의 생성·소비 관계와 상태
전이 검증을 이용하여 오염 관련 이력 추적을 다룬다 [4]. 이는
응용과 전이 구조 측면에서 가까운 기준선이다. 다만 그 구성에
서 공개되는 입력 문서 목록과 zkDPP에서 숨기려는 실제 부모
대응은 서로 다른 누출 조건이다. 비교할 때 비공개 문서 내용,
공개 연결 관계, 소유자 식별을 하나의 프라이버시 항목으로 합
치지 않아야 한다. zkDPP에서는 회계 상태를 한 번 소비하는
회계 증표·이전 증표와 외부 DPP의 지정 필드에 결합된 종단 제
품 주장을 구분하고, 이를 감사의 탐색 대상 및 현재 차단 가능
한 대상의 구분에 사용한다.

ZKVeil은 분산 식별자(decentralized identifier, DID)·자격
증명과 비공개 제품·금액 조건의 증명을 양자 거래의 주문·지불
·배송 상태에 결합한다 [6]. 거래의 제품 식별과 현재 세션 상태
검증은 본 연구에서도 고려해야 하는 기존 기능이다. 본 연구의
비교 대상은 그 위에 정의할 다단계 객체 이력의 복원과 전체 후
속 미소비 집합의 대응 절차이며, 개별 세션의 정상 완료와 이를
구분한다.

\subsection{비공개 거래의 추적과 승인된 공개}
GGM은 비공개 지급의 추적 정보를 증명 관계에 결합하고, 추
적 대상 입력이 소비되었을 때 후속 출력으로 추적을 이어가는
구성을 제시한다 [7]. 실제 출력과 연결된 암호화된 추적 정보
또는 후속 추적 자체를 본 연구의 새로운 원시기술로 볼 수 없
는 이유이다. Azeroth 역시 감사자가 거래의 감사 정보를 복호
화할 수 있도록 하고 그 정보와 실제 거래의 결합을 다룬다 [8].
이 기준선을 한 거래만 열 수 있는 감사 권한으로 제한하여 비교
하지 않는다.

Lee 등의 Token-Guided Flow Tracing은 방향별 추적 토큰
과 에포크 키를 이용하여 승인된 거래 흐름·방향·기간 안에서
감사 정보를 열도록 구성한다 [9]. 토큰 관리자와 감사자의 승
인 범위 밖 비공모를 가정하며, 수락된 노트의 소비 관계와 사용
자별 연속 거래를 따라 추적한다. zkDPP는 유효한 감사 데이터
키를 확보한 신뢰된 감사자가 같은 키 영역의 기록을 복호화하
는 모델이므로 이와 같은 암호학적 범위 제한을 제공하지 않는
다. 감사의 공개 범위와 그 결과를 사용 차단으로 연결하는 절

% Source PDF page 3, column 2
차는 서로 다른 비교 항목이다.

Derecho는 증명 동반 공개(proof-carrying disclosure)로
허용된 출처에 관한 증명을 거래의 실제 입력을 따라 재귀적으
로 전파한다 [10]. 이는 출처 조건의 증명과 거래 관계의 결합
에 관한 선행기술이며, 기존에 만들어진 증명을 사후 정책 변경
에 따라 취소하는 데 따르는 한계도 논의한다. zkDPP는 문제
시작점을 나중에 지정한 후 수락된 객체 그래프를 감사자가 복
원하는 절차를 다룬다. 기존 제품 주장을 취소하거나 과거 증
명을 소급하여 무효화하지 않고, 현재 미소비 객체의 상태를 이
후 소비 시점에 검사한다. 이 차이는 비교하는 보장의 범위이지
Derecho의 일반적 표현력에 대한 제한이 아니다.

\subsection{동결·리콜과 집행 대상의 차이}
CAP는 비공개 자산에 대해 정책별 입력·출력 열람과 동결 권한
을 정의한다 [5]. 열람용 메모(viewer memo)와 거래의 결합,
소유자와 동결 권한자가 같은 소비 식별자를 계산할 수 있는 구
조, 동결 권한자에 의한 레코드 소비·재발행은 zkDPP와 직접
비교해야 하는 요소이다. 열람 범위는 정책과 조건에 따르며,
동결은 하나 이상의 레코드를 동결 상태로 재발행하는 방식이
다. zkDPP의 현재 식별자에 대한 상태 변경과 구체 전이는 다
르지만, 비공개 거래의 열람과 사용 차단을 결합한다는 일반적
기능은 이미 CAP에 존재한다. CAP의 동결 정책은 전체 레코드
정보의 열람을 요구하며, 열람·동결 정책은 설정된 거래량 임곗
값의 적용을 받는다. 여러 입력을 한 동결 트랜잭션에서 소비·
재발행할 수 있다는 점도 비교에 포함해야 한다. zkDPP의 미래
식별자 사전 암호화와 CAP의 소유자·동결 권한자 간 식별자 계
산은 집행 정보를 얻는 서로 다른 구성이다. 전자는 출력 생성
자의 비밀 지식과 이전 증표 전달을 요구하며, 후자보다 자동으
로 적은 공개나 비용을 제공한다고 볼 수 없다.

ERC-20R은 공개 거래 이력에서 후속 금액을 추적하고 하나
의 원자적 동결 트랜잭션으로 회수 대상 금액을 확보하는 절차
를 제안한다 [11]. 이는 공개 계정과 금액 책임을 기준으로 하며,
혼합·분기된 공급망 객체의 모든 후속 출력에서 미소비 집합을
정하는 zkDPP의 규칙과 동일하지 않다. 특히 여러 상태 변경을
확인하는 zkDPP 실행에는 소비 경합과 부분 완료가 있으므로,
ERC-20R의 원자성을 그대로 보장한다고 서술하지 않는다.

ACTMeds는 영지식증명으로 리콜 사유를 검증하고 의약품
배치가 리콜 상태가 된 뒤의 전이를 제한하는 설계를 기술한
다 [12]. 공개된 설계 설명을 기준으로 한 비교이며, 본 연구에
서 그 실행 코드의 적합성을 검증한 것은 아니다. 여기서 배치
상태에 의한 리콜과 감사로 복원한 전체 후속 회계 증표·이전 증
표의 식별은 구분해야 한다. Circularise의 공개 특허는 비공개
의존 관계와 정보 유효성 유지에 대해 만료·취소·기록 동결 등
의 대안을 논의한다 [13]. 그중 기록 동결은 기록의 추가 갱신을
막는 의미이며, 디지털 자산의 소비 차단과 동일하지 않다. 특
허의 여러 대안을 하나의 구현이 모두 제공한다는 결과로 해석
하지 않는다.

\subsection{범용 구성과 본 연구의 기여 범위}
ZEXE와 VeriZexe는 비공개 상태의 생성·소비 규칙을 표현하
는 범용 기반을 제공한다 [14, 15]. zkDPP의 전이, 감사 암호문
결합, 상태 검사도 이러한 기반 위에서 구성할 수 있음을 인정한
다. 공정한 조합 기준선에는 동일한 감사자 신뢰, 동일한 키 접
근 범위와 공개 누출, 같은 회계·생애주기·완료 조건을 부여해

% Source PDF page 4, column 1
템이 그대로 구현해 놓은 결과를 가리키는 것은 아니다.

Table 1는 이 기준으로 가장 가까운 기능을 요약한다. 각 행
은 기능의 유무를 단정하는 목록이 아니라 집행 대상과 가정의
차이를 보여 준다. zkDPP가 제시하는 것은 생산·유통 객체의
실제 전이, 승인된 이력 복원, 미소비 집합의 현재 사용 차단을
함께 명세한 DPP 구성과 그 완료 조건의 분석이다. 개별 기법
의 최초성, 다른 구성으로의 표현 불가능성 또는 측정하지 않은
효율 우월성은 주장하지 않는다.

\subsection{동일 기능의 조합 기준선과 구성 선택}
조합 기준선 \(C^{\star }\)는 범용 비공개 생성·소비 관계에 본 논문과 동
일한 회계 정책, 감사 키 접근 및 공개 의미를 부여한 분석용 구
성이다. 그 관계에 실제 부모·출력 식별자의 암호화를 넣고, 원
장이 타입별 소비 기록과 상태를 검사하며, 감사자가 같은 전체
대상과 완료 조건을 사용하면 IX절의 논증은 \(C^{\star }\)에도 적용된다.
따라서 동일 기능의 기준선과 zkDPP 사이에 기능적 우위를 설
정하지 않는다.

구체화 과정에서 남는 선택은 다음과 같다. 첫째, 일반적인
수신자 소유 출력에서는 송신자가 미래 소비 식별자를 계산하
지 못할 수 있다. zkDPP는 회계 증표를 미래 소유자가 생성하
도록 제한하고 조직 간 이전을 별도의 이전 증표로 분리한다.
둘째, 생성 참조와 집행 식별자를 구분하므로 감사 평문에 타입
·위치·미래 식별자를 결합한다. 셋째, 상태 변경을 여러 트랜잭
션으로 수행하므로 정확한 추적과 최종 전체 완료를 별도로 판
정한다. 이 선택에는 이전 증표 개방값의 비밀 전달, 미래 식별
자 암호문, 현재 상태 조회와 결과 확인이라는 의무가 따른다.

이 의무를 제거하면 구성이 단순해질 수 있지만 보장도 달라
진다. 부모만 복호화하는 구성은 현재 미소비 객체의 집행 식별
자를 별도로 확보해야 하며, 일부 성공만 보고하는 구성은 전체
대상 차단을 뜻하지 않는다. 반대로 \(C^{\star }\)에 같은 의무를 모두 넣
으면 같은 보장을 얻는다. 본 논문의 구성 결과는 이러한 대응
을 실제 생애주기·자료 관계로 명시한 데 있으며, X절은 그 관계

% Source PDF page 4, column 2


\begin{table}[htbp]\centering\small
    \setcounter{table}{0}\caption{Closest mechanisms and the scope of comparison}
    \begin{tabularx}{\linewidth}{@{}XXX@{}}\toprule
        \textbf{Baseline}                          & \textbf{Related mechanism}                                             & \textbf{Distinction for comparison with zkDPP}                                        \\ \midrule
        Udokwu--Craß [3]                           & Private-input sustainability predicates                                & Accepted lifecycle graph and live-object response                                     \\
        Babel et al. [2]                           & Linked credentials, private ownership transfer, fraction consumption   & Post hoc targets and current consumption checks                                       \\
        Sahai et al. [4]                           & Supply-chain consume/create rules and public input document references & Hidden actual parent correspondence and authorized recovery                           \\
        GGM [7]                                    & Proof-bound tracing data and expansion to descendant outputs           & Typed lifecycle, enforcement targets and completion                                   \\
        Azeroth [8]                                & Transaction-bound encrypted audit information                          & Explicit graph traversal and whole-target response                                    \\
        Lee et al. [9]                             & Flow-, direction- and epoch-scoped audit capability                    & Key-domain-wide auditor access; separate containment procedure                        \\
        CAP [5]                                    & Policy-controlled viewing and freezing; consume-and-reissue freeze     & Full-reveal/threshold conditions; shared spend ID versus encrypted future ID          \\
        ERC-20R [11]                               & Downstream amount tracing and atomic freezing on public history        & Account liability versus full descendant sets; atomic execution versus spending races \\
        ACTMeds [12]                               & Private recall-reason proofs; transitions restricted by batch status   & Batch-level restrictions versus all recovered live descendants                        \\
        ZEXE family and composed baseline [14, 15] & Generic private creation and consumption rules                         & Same trust, key access and leakage; same application relations                        \\
        \bottomrule\end{tabularx}\end{table}
